% !TEX root = ../PhD Thesis.tex
\chapter{SGs and Cancer}
\label{chap:ChapterV}

\section{Inflammation: A Double Agent in Plain Sight}

The previous chapter suggested that \glspl{sg} can modulate inflammatory signalling. Yet, their impact net impact on cancer remains unclear. Inflammation itself is a highly context-dependent process in oncology: it can be tumour-promoting, when chronic inflammatory signalling drives genomic instability and fosters a mutagenic microenvironment conducive to cancer evolution and therapy resistance; or tumour-suppressing, when inflammatory mediators recruit immune effector mechanisms capable of eliminating transformed cells\cite{turizo-smith_inflammation_2024}.
\\
\\
Importantly, \glspl{sg} are embedded within the broader integrated stress response, which is tightly coupled to multiple inflammatory signalling pathways. As such, their effects on tumour biology are likely not independent of inflammation, but may instead reflect specific modulation of stress and immune signalling networks. Following the approach taken throughout this thesis, we therefore revisited several studies directly implicating \glspl{sg} in cancer, with the aim of assessing whether their reported effects are similarly context-dependent.

\section{Materials and Methods}

\subsection{TCGA Expression Data}

\gls{rnaseq} data (in FASTQ format) generated by the TCGA Research Network were downloaded from the \gls{gdc} Data Portal\cite{network_cancer_nodate, heath_nci_2021}. Raw reads were trimmed using Trimmomatic (v0.39)\cite{bolger_trimmomatic_2014}. Trimmed reads were then aligned to the \gls{hg38} reference genome using \gls{star} (v2.7.9a)\cite{dobin_star_2013} with the following parameters:

\begin{lstlisting}
	--outFilterType BySJout
	--alignSJoverhangMin 8
	--alignSJDBoverhangMin 1
	--alignIntronMin 20
	--alignIntronMax 1000000
	--alignMatesGapMax 1000000
	--outFilterMultimapNmax 1
	--outFilterMismatchNmax 999
	--outFilterMismatchNoverLmax 0.04
\end{lstlisting}
\noindent
Aligned reads (\gls{bam} files) were subsequently annotated using the GENCODE v42 primary assembly\cite{frankish_gencode_2021} and quantified with featureCounts (v2.0.3)\cite{liao_featurecounts_2014}.
\\
\\
Read count data for each cancer type were pre-processed separately. We started by filtering out genes deemed uninformative due to their very low expression across samples: only genes with at least one CPM in at least 40\% of the samples were kept for analysis, akin to what we previously described in other large scale datasets\cite{schneider_voyager_2024}. Read counts for those kept genes were then normalised as described in Section 3.2.3.

\subsection{DGEA}

To model expression changes between conditions, we fitted linear models to the gene expression data, that accounted for sample type (primary tumour vs normal) and tissue of origin (breast, skeletal muscle, etc). Standard errors were then moderated using empirical Bayes shrinkage\cite{smyth_linear_nodate}, and \gls{deg}s were defined as indicated on the figures.

\subsection{GSEA}

\gls{gsea}\cite{subramanian_gene_2005} was performed using the fgsea (v1.20.0)\cite{korotkevich_fast_2016} R package. Genes were ranked by t-statistics from \gls{dgea}, ordered from most positive to most negative, and analysed against the Hallmark gene sets\cite{liberzon_molecular_2015}. The \gls{nes} was extracted to quantify the extent to which a given gene set is enriched at the extremes of the ranked gene list.

\subsection{Alternative Hypothesis Testing for SG Implication in Cancer}

To explore mechanisms beyond \gls{sg} assembly that could account for the phenotypes observed in \gls{ptgs2} knockout cells (which were used to abbrogate SG assembly in its respective study), specifically drug sensitivity, multiple computational resources were leveraged. \gls{ptgs2} perturbagens were identified using the \gls{cmap}\cite{subramanian_next_2017}. Drug sensitivity profiles were obtained from the \gls{gdsc}\cite{yang_genomics_2012} database. 

\subsection{SG Signature}

To define a machine learning-derived expression signature indicative of \gls{sg} presence, a discovery dataset was first constructed by integrating GSE138988\cite{matheny_transcriptome-wide_2019} and GSE119977\cite{matheny_rna_2021, matheny_corrigendum_2023}. This integration enabled comparison of stressed cells with and without \glspl{sg} while controlling for confounding variables including cell type, genotype, stressor, and study of origin. Batch effects arising from these variables were corrected prior to analysis following the methodology described previously in Chapter 4\cite{schneider_voyager_2024}. An elastic net regression was then applied with 1000 bootstrap iterations and a mixing parameter (\alpha) of 0.75, yielding a more stringent feature selection compared to a more commonly used $\alpha$ of 0.5\cite{zou_regularization_2005}. Genes selected in at least 900 of 1000 iterations were retained\cite{zou_regularization_2005}. The resulting signature comprised positively weighted genes (upregulated in \gls{sg}-positive cells) and negatively weighted genes (downregulated in \gls{sg}-positive cells), capturing the bidirectionality of the \gls{sg} transcriptomic response.
\\
\linebreak
\noindent
The \gls{sg} score for individual samples was computed using a non-parametric gene ranking approach, with the markeR R package\cite{martins-silva_exploring_2025}. Within each sample, all genes were ranked by expression level. The score was then calculated as the sum of expression ranks of signature genes, weighted by $+1$ or $-1$ for positively and negatively weighted genes respectively, and normalised by the total number of genes in the signature. The signature was validated on an independent dataset, GSE173953\cite{paget_stress_2023}, applying identical batch effect correction as described in Chapter 4\cite{schneider_voyager_2024}.

\subsection{Survival Analysis}

For survival analysis, \gls{gdc}-\gls{tcga} samples\cite{network_cancer_nodate, heath_nci_2021} were stratified into top and bottom quartiles based on the expression of SG-implicated genes (in section 6.3.1), or on \gls{sg} Score (in section 6.4.2). Kaplan-Meier survival curves were generated and statistical significance assessed using the log-rank test\cite{goel_understanding_2010}, corrected for multiple testing using \gls{bh}. Effect size was assessed in the form of \gls{hr}.
\\
\\
To evaluate the robustness of observed survival associations, 100 random signatures matching the original \gls{sg} signature in size and gene directionality\cite{martins-silva_exploring_2025} were generated by sampling gene names from the \gls{tcga} \gls{laml} dataset\cite{network_cancer_nodate, heath_nci_2021}, selected for its large number of genes passing expression filters. Survival analyses were repeated for each random signature, and false discovery rates were computed as the proportion of random signatures yielding log-rank p-values below those obtained from the original \gls{sg} signature.


\subsection{\textit{In Silico} Resources}

All analyses, unless otherwise specified, were performed on the workstation \textit{lobito} at the \gls{gimm} data center, which is equipped with 19 \gls{tb} of \gls{ssd} storage, 768 \gls{gb} of RAM, 72 \gls{cpu} threads (\gls{intel} Xeon
Gold 6254 @ 3.10 \gls{ghz}), and an integrated Matrox G200eW3 \gls{gpu}. \gls{bash} was used to run FASTQC\cite{andrews_fastqc_2010}, MultiQC\cite{ewels_multiqc_2016}, STAR\cite{dobin_star_2013}, and featureCounts\cite{liao_featurecounts_2014}. Downstream analyses were conducted in R (v4.1.2) using RStudio Server (v1.4.1717). 

\section{The Problem with Single-Gene Approaches in SG Studies}
\subsection{In Metastasis}
Given the established link between inflammation and cancer\cite{turizo-smith_inflammation_2024}, the involvement of \glspl{sg} in inflammatory regulation may help explain their reported association with tumour progression. Namely, \glspl{sg} have been proposed to enhance the metastatic potential of sarcoma cells\cite{somasekharan_yb-1_2015}, based on comparisons between \gls{sg}-competent and \gls{sg}-deficient cells. In this study, \gls{sg} ablation was achieved via \gls{g3bp1} \gls{ko}, a protein with functions extending beyond \gls{sg} assembly, which may have confounded the observed phenotype. Furthermore, \gls{g3bp1} \gls{ko} cells retain the capacity to form \glspl{sg} under certain stressors, such as hyperosmotic stress\cite{matheny_rna_2021, matheny_corrigendum_2023}, a condition potentially encountered in the tumour microenvironment.
\\
\\
A recent study has similarly implicated \glspl{sg} as tumour promoting in pancreatic cancer, using comparisons between \gls{sg}-competent and \gls{sg}-deficient cells\cite{fonteneau_stress_2022}. In this case, rather than generating a full \gls{g3bp1} \gls{ko}, the authors specifically disrupted its \gls{ntf2l} domain required for \gls{sg} assembly, thereby aiming to more selectively isolate \gls{sg}-related functions. Still, it is possible that additional functions of \gls{g3bp1} independent of \gls{sg} formation also depend on this domain and may contribute to the observed phenotype. To strengthen their conclusions, the study additionally targeted a secondary \gls{sg}-nucleating protein, \gls{tiar}. While this increases robustness, both proteins are \glspl{rbp} with potentially overlapping functions, raising the possibility that the observed effects reflect shared \gls{rbp} activities rather than \gls{sg}-specific mechanisms.
\\
\linebreak
\noindent
Multiple genes, including \textit{\gls{g3bp1}}\cite{tourriere_rasgap-associated_2023}, \textit{\gls{ago2}}\cite{leung_quantitative_2006}, \textit{\gls{brf1}}\cite{stoecklin_mk2induced_2004}, \textit{\gls{caprin1}}\cite{solomon_distinct_2007}, \textit{\gls{cpeb1}}\cite{wilczynska_translational_2005}, \textit{\gls{ddx6}}\cite{wilczynska_translational_2005}, \textit{\gls{disc1}}\cite{ogawa_functional_2005}, \textit{\gls{eif2a}}\cite{kedersha_rna-binding_1999}, \textit{\gls{fastk}}\cite{kedersha_stress_2005}, \textit{\gls{fmr1}}\cite{mazroui_trapping_2002}, \textit{\gls{fxr1}}\cite{mazroui_trapping_2002}, \textit{\gls{g3bp2}}\cite{tourriere_rasgap-associated_2023}, \textit{\gls{l1td1}}\cite{goodier_line-1_2007}, \textit{\gls{samd4a}}\cite{baez_mammalian_2005}, \textit{\gls{smn1}}\cite{hua_survival_2004}, \textit{\gls{tia1}}\cite{gilks_stress_2004}, \textit{\gls{tiar}}\cite{gilks_stress_2004}, and \textit{\gls{ttp}}\cite{stoecklin_mk2induced_2004}, are known to nucleate \glspl{sg} spontaneously upon overexpression\cite{anderson_stress_2008}. We hypothesised that if \gls{sg}s directly drive tumour aggressiveness, high expression of these \gls{sg}-nucleating genes should correlate with poorer patient survival.
\\
\\
Analysis of the GDC-TCGA\cite{network_cancer_nodate, heath_nci_2021} sarcoma dataset (the cancer type assessed in the metastatic study described above), considering only genes present after expression filtration, revealed that of the eight genes that survived filtration, only \textit{AGO2} expression was significantly associated with patient survival after \gls{bh} correction for multiple testing (HR = 2.43, p\textsubscript{BH} = 0.0178; Figure 6.1). The hazard ratio reflects the risk associated with high \textit{AGO2} expression relative to low, with higher expression linked to worse survival outcomes. While this direction of effect is nominally consistent with a pro-tumourigenic role of \gls{sg} assembly, in that higher expression of an \gls{sg}-nucleating gene correlates with poorer prognosis, it is difficult to attribute this association to \gls{sg} activity specifically, given the well-established \gls{sg}-independent functions of \textit{AGO2} as a core component of \gls{risc}, a gene silencing complex, whose dysregulation has been independently linked to cancer progression\cite{iwakawa2022life}. Notably, \textit{G3BP1}, the primary \gls{sg} nucleator implicated in the sarcoma metastasis study, showed no significant association with survival. Given that only one of the eight tested genes reached significance, and that its prognostic effect is more parsimoniously explained by its broader role in RNA silencing, this analysis does not provide strong support for a direct role of \gls{sg} assembly in sarcoma aggressiveness. More broadly, even for genes with a clear role in \gls{sg} nucleation, their expression level is an imperfect proxy for \gls{sg} activity, as the functional consequences of \gls{sg} assembly may be buffered or masked by the pleiotropic roles of these proteins in other cellular processes. Ultimately, without a direct measure of \gls{sg} formation (or at least a transcriptomic signature capable of serving as a reliable proxy), the prognostic relevance of \glspl{sg} in sarcoma (and other cancer types) remains speculative.


\begin{figure}[H]
	\centering
	\includegraphics[width=1\linewidth]{"images/Paper Images/Figure10_2"}
	\caption[Overall Survival Analysis Based on the Expression of Genes Encoding SG-Related Proteins]{\textbf{Overall survival analysis based on the expression of genes encoding SG-related proteins.} Kaplan–Meier plots showing proportion of overall survival (Y-Axis) per day (X-Axis) of patients from GDC-TCGA's \gls{sarc} dataset\cite{network_cancer_nodate, heath_nci_2021}, stratified by expression of the indicated genes, as plot titles. Patients were grouped into high and low expression cohorts based on the bottom and top quartiles expression values. \gls{bh}-adjusted Log-rank p-values, and \gls{hr}, are shown in the top right corner of each plot. Genes with an adjusted Log-rank p-value $<$0.05 are shown in red.}
	\label{fig:Figure8.1}
\end{figure}


\subsection{In Chemoresistance}
\glspl{sg} have also been proposed to promote chemotherapy resistance\cite{grabocka_mutant_2016}. In this model, \gls{kras}-mutant cells enhance \gls{sg} assembly paracrinally via \textit{\gls{ptgs2}} upregulation\cite{grabocka_mutant_2016}. Again, \textit{\gls{ptgs2}} has established \gls{sg}-independent roles, and the observed resistance phenotype may not be solely attributable to \gls{sg} assembly. {\gls{kras}-mutant cells exhibit constitutive activation of the \gls{egfr} pathway\cite{huang_kras_2021}. As \gls{kras} functions downstream of \gls{egfr}, its mutation predicts poor response to \gls{egfr} inhibitors but improved response to \gls{mek} inhibitors, which act further downstream\cite{hondo_mek_2023} (Figure 6.2A, B). \gls{ptgs2} \gls{kd} or inhibition produces expression changes that correlate with \textit{\gls{egfr}} and \textit{\gls{erbb3}} \gls{kd}, another member of the same receptor family, suggesting that \gls{ptgs2} may operate within the same signalling pathway (Figure 6.2C). If \gls{ptgs2} acts downstream of \gls{kras}, its inhibition may functionally mimic \gls{mek} inhibition, targeting a key oncogenic driver and thereby accounting for the observed increase in treatment sensitivity. Although \gls{ptgs2} is known to be regulated by the \gls{egfr} pathway\cite{wang_crosstalk_2011}, this proposed mechanism remains to be experimentally validated.

\begin{figure}[H]
	\centering
	\includegraphics[width=1\linewidth]{"images/Paper Images/Figure11"}
	\caption[KRAS mutation and associated drug sensitivities and resistances]{}
	\label{fig:Figure8.2}
\end{figure}

\clearpage


\begin{figure}[t]
	\captionsetup{labelformat=adja-page}
	\ContinuedFloat
	\caption[]{\textbf{KRAS mutation and associated drug sensitivities and resistances.} (A) Volcano plot of \gls{ic50} (X-Axis -  IC50 Effect) and corresponding significance, given by MANOVA (Y-Axis – log10(p-value)) in \gls{kras} mutant cells\cite{yang_genomics_2012}; drugs/chemicals as circles, with drugs with the most differential \gls{ic50} effect labelled and coloured in green (if with increased sensitivity) or in red (if with increased resistance). Green circle around drugs that target \gls{mek}, and red circles around drugs that target \gls{egfr}. Analysis and images are obtained from \gls{gdsc}\cite{yang_genomics_2012}, using the \gls{gdsc}2 dataset. (B) Simplified schematic of the \gls{egfr} signalling pathway. (Left) In \gls{wt} cells, \gls{egfr} activation triggers \gls{kras} activation, which in turn activates \gls{mek} and downstream signaling. (Middle) In \gls{kras}-mutant cells, inhibition of \gls{egfr} with targeted drugs (red drugs, see (A)) does not block signaling, as mutant \gls{kras} activates \gls{mek} independently. This explains resistance of \gls{kras}-mutant cells to anti-\gls{egfr} therapy. (Right) In \gls{kras}-mutant cells, inhibition of \gls{mek} with targeted drugs (green drugs, see (A)) blocks downstream signaling, rendering these cells sensitive to anti-\gls{mek} therapy. (C) Table representing alterations with highest correlation to \gls{ptgs2} \gls{kd}. Values represent \gls{cmap} connectivity scores, which reflect the degree of similarity between gene expression signatures. \glspl{kd} of \gls{egfr} and \gls{erbb3} are highlighted in red. Analysis and visualisation were performed and adapted from the \gls{cmap} platform\cite{subramanian_next_2017}.}
\end{figure}
\noindent

\section{The SG Score}

\subsection{Constructing a SG Transcriptomic Signature}

The results above highlight the limitations of single-gene analyses in isolating \gls{sg}-specific effects, given the overlapping roles of individual \gls{sg} components in broader stress-related (or even unrelated) pathways. To address this, a custom \gls{sg} transcriptomic signature was developed. Datasets  GSE138988\cite{matheny_transcriptome-wide_2019} and GSE119977\cite{matheny_rna_2021, matheny_corrigendum_2023} were integrated, as this combination provided stressed samples with and without \glspl{sg}, enabling batch effect correction for cell genotype, stressor type, and study of origin. An elastic net regression model with bootstrapping was then applied to derive the signature (see 6.2.6). Genes retained in more than 90\% of bootstrap iterations were used to construct the final 41-gene signature, hereafter referred to as the \gls{sg} Score. The signature comprised genes expected to be upregulated upon \gls{sg} assembly, alongside a larger set of genes expected to be downregulated, which often resulted in overall negative \gls{sg} Scores. Both the signature and scoring approach were validated on an independent dataset (GSE173953\cite{paget_stress_2023}), which, despite originating from a single study, required batch correction due to processing across three distinct batches (as indicated by submission dates on \gls{geo}). The \gls{sg} Score successfully discriminated \gls{sg}-positive from \gls{sg}-negative samples in this dataset (Figure 6.3).
\\
\\
Importantly, the \gls{sg} Score should be interpreted as a relative, rather than absolute, measure of \gls{sg} presence. Higher scores indicate stronger support for \gls{sg} assembly within a given dataset, but do not define a universal threshold applicable across studies. For example, \gls{sg}-positive samples in the discovery dataset exhibit scores of approximately -12.5, whereas non-\gls{sg} samples in the validation dataset can display similar values (Figure 6.3). This discrepancy likely reflects differences in cell type, stress conditions, and residual inter-dataset variation, even after batch correction, which shift the baseline expression of the signature genes and consequently the baseline score values. The \gls{sg} Score should therefore be used to rank or compare samples in the same conditions, rather than to make absolute comparisons across different conditions.

\begin{figure}[H]
	\centering
	\includegraphics[width=1\linewidth]{"images/Paper Images/TESE_Fig8.3"}
	\caption[Development and Validation of the SG-Score]{\textbf{Development and Validation of the SG-Score.} Violin plots of \gls{sg} signature scores, calculated through a ranking metric\cite{martins-silva_exploring_2025}, in cells without and with \glspl{sg}. \gls{wc} samples are coloured in accordance to stress type stimuli. The discovery dataset (A) refers to GSE138988\cite{matheny_transcriptome-wide_2019} and GSE119977\cite{matheny_rna_2021, matheny_corrigendum_2023}  whereas as the validation dataset (B) refers to GSE173953\cite{paget_stress_2023}. Effect Size (Cohen's d) and corresponding statistical significance, given by Welch's t-test (p-value), are shown above plots.}
	\label{fig:Figure8.3}
\end{figure}
\noindent

\subsection{SGs as Survival Biomarkers}

The prognostic performance of the \gls{sg} signature was evaluated across \gls{gdc}-\gls{tcga} datasets\cite{network_cancer_nodate, heath_nci_2021}, encompassing the following cancer types: \gls{laml}, \gls{acc}, \gls{chol}, \gls{blca}, \gls{brca}, \gls{cesc}, \gls{coad}, \gls{ucec}, \gls{esca}, \gls{gbm}, \gls{hnsc}, \gls{kich}, \gls{kirc}, \gls{kirp}, \gls{dlbc}, \gls{lihc}, \gls{lgg}, \gls{luad}, \gls{lusc}, \gls{skcm}, \gls{meso}, \gls{uvm}, \gls{ov}, \gls{paad}, \gls{pcpg}, \gls{prad}, \gls{read}, \gls{sarc}, \gls{stad}, \gls{tgct}, \gls{thym}, \gls{thca}, and \gls{ucs}.
\\
\\
Given that cancer cells frequently encounter conditions known to promote \gls{sg} formation\cite{anderson_stress_2015, szaflarski_vinca_2016, moeller_radiation_2004}, we hypothesised that the \gls{sg} Score would be elevated in tumour samples relative to matched normal tissue\footnote{It is important to note that, in \gls{tcga}, normal tissue samples are obtained from close proximity to the primary tumour (around 2 cm)\cite{network_cancer_nodate, heath_nci_2021}. It is plausible that this ``normal" tissue is under increased stress and can thus exhibit \gls{sg}s as well. Nonetheless, as this stress is theoretically lower than what is expected on tumour cells, this observation is simply noted here, as a potential confounder.}. Among the 17 \gls{tcga} cancer types with at least five normal and five tumour samples, only three, \gls{blca}, \gls{lihc}, \gls{prad}, exhibited significantly increased \gls{sg} Scores in tumour samples ($Welch's t-test, adj.p<0.05$). In contrast, seven cancer types, \gls{esca}, \gls{hnsc}, \gls{kirc}, \gls{kirp}, \gls{luad}, \gls{lusc}, and \gls{stad}, displayed significantly lower \gls{sg} Scores in tumours relative to normal tissue (p-value$<0.05$), with \gls{esca} and \gls{stad} showing particularly large effect sizes (Cohen's d $< -1$) (Figure 6.4).

\begin{figure}[H]
	\centering
	\includegraphics[width=1\linewidth]{"images/Paper Images/Sup10"}
	\caption[SG score across GDC-TCGA datasets]{}
	\label{fig:Figure8.4}
\end{figure}

\clearpage


\begin{figure}[t]
	\captionsetup{labelformat=adja-page}
	\ContinuedFloat
	\caption[]{\textbf{SG Score across GDC-TCGA datasets.} Violin plots of \gls{sg} signature scores, calculated through a ranking metric\cite{martins-silva_exploring_2025}, in tumour (blue) and normal tissue (orange) samples from several GDC-TCGA datasets\cite{network_cancer_nodate, heath_nci_2021}. Effect Size (Cohen's d) and corresponding statistical significance given by Welch's t-test (adjusted p-value) are shown above plots. Tumour types with an adjusted p-value $<$0.05 are shown in red. Only GDC-TCGA\cite{network_cancer_nodate, heath_nci_2021} datasets with at least 5 normal and 5 tumor samples were considered for this analysis.}
\end{figure}
\noindent
To assess prognostic relevance, survival analyses were performed comparing the top and bottom quartiles of the \gls{sg} Score. Although the score is continuous and should be interpreted in a relative manner, with higher values indicating greater support for \gls{sg} presence rather than an absolute state, the score has inherently limited precision and \glspl{sg} may not impose a sufficiently strong transcriptional phenotype to produce detectable differences across the full score distribution. We therefore focused on the extreme quartiles, between which any hypothesised biological differences are expected to be strongest, in order to maximise the contrast between samples with clearly low and clearly high inferred \gls{sg} activity. The underlying hypothesis remains that more aggressive or therapy-resistant tumours would exhibit higher scores and thus poorer overall survival\cite{somasekharan_yb-1_2015, grabocka_mutant_2016}.
\\
\\
As random gene signatures have been shown to exhibit spurious prognostic value in cancer\cite{venet_most_2011, shimoni_association_2018}, the \gls{sg} signature was benchmarked against 100 randomly generated signatures of matched size and directionality\cite{martins-silva_exploring_2025}. The \gls{sg} signature achieved an adjusted p.value and \gls{fpr} < 0.05 in two cancer types: \gls{acc} and \gls{meso}, suggesting that \gls{sg} presence may relate to tumour aggressiveness in these contexts (Figure 6.5). In both cancer types, higher \gls{sg} Scores were consistently associated with improved patient survival, which is contrary to the hypothesis that \gls{sg} assembly promotes tumour aggressiveness. One possible interpretation is that \gls{sg} formation in these tumour types reflects an intact stress response capacity, which may in turn correlate with a less aggressive tumour phenotype or greater sensitivity to therapy. Alternatively, the prognostic signal may be driven by \gls{sg}-independent functions of the genes comprising the signature, rather than by \gls{sg} activity per se. These possibilities are not mutually exclusive, and distinguishing between them would require direct experimental assessment of \gls{sg} formation in these tumour contexts.

\begin{figure}[H]
	\centering
	\includegraphics[width=1\linewidth]{"images/Paper Images/Sup11_2"}
	\caption[Survival Stratification with SG Score, across GDC-TCGA datasets]{}
	\label{fig:Figure8.5}
\end{figure}

\clearpage


\begin{figure}[t]
	\captionsetup{labelformat=adja-page}
	\ContinuedFloat
	\caption[]{\textbf{Survival stratification with SG Score, across GDC-TCGA datasets.} Kaplan–Meier plots showing proportion of overall survival (Y-Axis) per days (X-Axis) of patients of GDC-TCGA datasets\cite{network_cancer_nodate, heath_nci_2021}, stratified by SG score, with dataset identification as plot title. Patients were grouped into high and low expression cohorts based on bottom and upper quartile SG Score values. BH-adjusted Log-rank p-value, HR, and FPR compared to 100 random signatures is shown on bottom right of each plot; datasets with FPR $>$ 0.05 are highlighted in red.}
\end{figure}
\noindent

\subsection{Mechanistic Basis of SG Prognostic Associations}

To investigate why prognostic significance was restricted to \gls{acc} and \gls{meso}, gene expression in these two cancer types was compared against the remaining GDC-TCGA tumour cohorts\cite{network_cancer_nodate, heath_nci_2021}, controlling for tissue of origin. This revealed strong upregulation of \textit{\gls{mamstr}} (Figure 6.6A), a \gls{mef2} coactivator previously linked to \gls{emt} and improved survival\cite{creemers2006coactivation, pon2015mef2}, alongside downregulation of proliferation-associated Hallmarks and upregulation of \gls{emt} and inflammation-related programmes (Figure 6.6B). The association between elevated inflammatory signalling and improved survival in high \gls{sg} Score samples was unexpected, given the proposed role of \glspl{sg} in attenuating inflammation\cite{paget_stress_2023}. One possibility is that heightened inflammatory signalling in these tumour types potentiates anti-tumour immunity, thereby restricting tumour growth\cite{nishida_role_2025}, and that \gls{sg} assembly reflects or accompanies this state.
\\
\\
To further characterise the transcriptional differences associated with \gls{sg} Score within these cancer types, samples in the top and bottom \gls{sg} Score quartiles were compared by \gls{dgea}, controlling for tissue of origin. This identified upregulation of \textit{\gls{hk2}} (Figure 6.6C), which associates with mitochondrial outer membrane permeabilisation through binding of \gls{vdac}\cite{haloi_structural_2021}. \gls{gsea} revealed a partial inversion relative to the basal state (Figure 6.6D): tumours with higher \gls{sg} Scores showed enrichment of proliferation-associated programmes that were globally downregulated in these cancer types compared to other tumours. Interferon-\alpha and interferon-\gamma responses remained elevated in high \gls{sg} Score samples, which might appear inconsistent with an anti-inflammatory role for \glspl{sg}. However, this may reflect a comparison between tumours operating under differing levels of cellular stress, in which \glspl{sg} attenuate but do not abolish inflammatory signalling\footnote{Additionally, the relationship between a molecular feature and its prognostic value is not necessarily linear; features that are globally reduced in tumours may still stratify outcomes within the tumour population. This behaviour has been described for tumour suppressor genes, which, despite being globally downregulated in cancer, can associate with differential prognosis when relatively preserved\cite{datta_tumor_2020, noauthor_sparc-p53_2020}.}.


\begin{figure}[H]
	\centering
	\includegraphics[width=1\linewidth]{"images/Paper Images/TESE_Fig8.6"}
	\caption[Mechanistic Insights on SGs as Prognostic Markers]{\textbf{Mechanistic insights on SGs as prognostic markers.} (A) Volcano plot of differential gene expression (X-Axis -  log2 Fold-Change;Y-Axis – B-Statistic) between \gls{acc}, and \gls{meso} tumour samples, and  remaining \gls{gdc}-\gls{tcga} tumour cohorts\cite{network_cancer_nodate, heath_nci_2021}, corrected for dataset of origin; genes as circles; most differentially expressed genes are coloured red (if upregulated, Log2FC$>$1.5 with B-Statistic$>$300) or blue (if downregulated, Log2FC$< -1.75$ with B-Statistic$>$300); (B)  \gls{gsea} \gls{nes} of Hallmark gene sets, ranked by t-statistic, between \gls{acc}, and \gls{meso} tumour samples, and remaining \gls{gdc}-\gls{tcga} cohorts\cite{network_cancer_nodate, heath_nci_2021}, corrected for dataset of origin; (C) Volcano plot of differential gene expresion (X-Axis -  Log2 Fold-Change; Y-Axis – B-Statistic) between top 25\% and  bottom 25\% of SG-Score tumour samples in \gls{acc}, and \gls{meso}, corrected for dataset of origin; genes as circles; most differentially expressed genes are coloured red (if upregulated, log2FC$>$1 with B-Statistic$>$0) or blue (if downregulated, Log2FC$< -1.5$ with B-Statistic$>$0); (D)  \gls{gsea} \gls{nes} of Hallmark gene sets, ranked by t-statistic, between top 25\% and bottom 25\% of SG-Score tumour samples in \gls{acc}, and \gls{meso}, corrected for dataset of origin.}
	\label{fig:Figure8.6}
\end{figure}
\noindent

